\section{Identification}
\label{sec:identification}

The introduction posed two questions. Where in the cross-section does the literacy mechanism become visible? And what microfoundation does that location identify? This section describes the empirical design that answers them.

\subsection{What the research question requires}

The hypothesis is that state-level investor literacy moderates the momentum--idiosyncratic-volatility (IV) anomaly, with the dampening larger in literate states. Testing this hypothesis means measuring how the within-stock relationship between returns and the $\mathrm{mom}\times\mathrm{IV}$ interaction varies as a function of the headquarters-state literacy of each stock. Three features of the data dictate the design.

First, literacy is measured at the state-wave level: the FINRA National Financial Capability Study fields the Big-Five literacy questions every three years and aggregates responses to the state level, so the within-state literacy variation is coarse (five wave values over the 2009--2023 sample) and the cross-state variation is the primary source of identifying variation. Anomaly returns are measured at the firm-month level. Bridging these scales requires the residual variation that identifies the moderation to be cross-firm \emph{within} state-month cells --- not cross-state, which would conflate the literacy channel with any other state-level determinant of returns, and not cross-time, which would conflate it with aggregate market dynamics.

Second, the channel transmits through local-investor community ties, which take time to form and to decay. For firms that relocate headquarters during the sample, the effective local-investor base is in transition: neither the origin nor the destination state correctly characterizes the marginal local retail investor. Including relocating firms under a static-snapshot headquarters assignment contaminates the literacy variable for the 14 percent of firm-months whose static state differs from their then-current state. A clean design therefore restricts to firms with stable headquarters across the panel window and reports the static-snapshot vs.\ point-in-time correction as a falsification check.

Third, the prediction is about where in the cross-section the moderation operates --- specifically, about how the moderation varies with firm ownership structure. The natural conditioning variable is institutional ownership share, but the conditioning is conceptually distinct from the headline interaction. A clean design therefore stratifies the regression by within-stock institutional-ownership tercile and reports the moderation coefficient separately within each tercile, rather than collapsing across IO levels.

\subsection{The baseline specification}

The headline specification implements these requirements as a two-way fixed-effects panel regression:

\begin{equation}
\label{eq:headline}
\begin{split}
r_{i,s,t} = {}& \alpha_{s} + \tau_{t} + \gamma \,\bigl(\mathrm{mom}_{i,t}\cdot\mathrm{IV}_{i,t}\cdot\mathtt{literacy\_z}_{s(i,t),t}\bigr) \\
& + \delta_{1}\,(\mathrm{mom}_{i,t}\cdot\mathrm{IV}_{i,t}) + \delta_{2}\,(\mathrm{mom}_{i,t}\cdot\mathtt{literacy\_z}_{s(i,t),t}) + \delta_{3}\,(\mathrm{IV}_{i,t}\cdot\mathtt{literacy\_z}_{s(i,t),t}) \\
& + \boldsymbol{\beta}'\mathbf{X}_{i,s,t} + \varepsilon_{i,s,t}.
\end{split}
\end{equation}

The dependent variable $r_{i,s,t}$ is the month-$t$ return of firm $i$ headquartered in state $s(i,t)$. State fixed effects $\alpha_{s}$ absorb time-invariant state characteristics --- demographics, industrial mix, the entire stock of state-specific factors that would otherwise contaminate the literacy variable. Month fixed effects $\tau_{t}$ absorb aggregate time-varying shocks. The three two-way interactions enter as standard controls; the three-way coefficient $\gamma$ is the object of interest. The state-assignment function $s(i,t)$ is the static Compustat snapshot under the headline merge and the point-in-time EDGAR 10-K-header reassignment under the falsification check. On the stable-headquarters subsample, the two coincide. The discrete-tercile rendering of Equation~\eqref{eq:headline} re-estimates the coefficient within each within-stock institutional-ownership tercile.

\subsection{Adding size as a moderator: why and how}

Mid-IO firms in the stable-HQ panel are concentrated in mid-cap market-equity quintiles: median firm market capitalization is roughly five to ten times higher in mid-IO than in low-IO. A reader could reasonably ask whether the mid-IO three-way is a size moderation operating through the same anomaly, captured by the IO discretization because mid-IO and mid-cap firms substantially overlap. The size-controlled specification rules this out by entering size as a parallel moderator on the same anomaly:

\begin{equation}
\label{eq:size_ctrl}
\begin{split}
r_{i,s,t} = {}& \alpha_{s} + \tau_{t} + \gamma \,\bigl(\mathrm{mom}_{i,t}\cdot\mathrm{IV}_{i,t}\cdot\mathtt{literacy\_z}_{s(i,t),t}\bigr) \\
& + \zeta \,\bigl(\mathrm{mom}_{i,t}\cdot\mathrm{IV}_{i,t}\cdot\mathtt{size\_z}_{i,t}\bigr) \\
& + \text{lower-order literacy interactions} + \text{lower-order size interactions} + \varepsilon_{i,s,t},
\end{split}
\end{equation}

where $\mathtt{size\_z}_{i,t}$ is the within-month $z$-score of $\log(\mathrm{ME}_{i,t})$. Size enters with the same standardization and the same interaction structure as literacy. The coefficient $\gamma$ in Equation~\eqref{eq:size_ctrl} is identified holding size moderation of the same anomaly fixed. If the literacy moderation is mechanically conflated with size moderation through the mid-IO sample's mid-cap concentration, $\gamma$ should attenuate to insignificance under the size-controlled specification; if literacy is a distinct moderator operating through a different channel, $\gamma$ should survive. The results below report both Equation~\eqref{eq:headline} (baseline) and Equation~\eqref{eq:size_ctrl} (size-controlled). Equation~\eqref{eq:size_ctrl} is the preferred specification throughout the paper.

\subsection{Identifying assumptions}

For the three-way coefficient $\gamma$ to be interpretable as a literacy moderation rather than as something else, three assumptions must hold.

\paragraph{Conditional unconfoundedness for the literacy moderator.} Conditional on state and month fixed effects, the three two-way interactions, and the size moderator in Equation~\eqref{eq:size_ctrl}, residual variation in state-aggregate literacy is uncorrelated with omitted determinants of returns that interact with $\mathrm{mom}\times\mathrm{IV}$ in the same functional form. The paper addresses this assumption four ways. First, the Fama-French 12-industry block enters the robustness battery as an additional control. Second, a state college-attainment horse-race within the headline tercile returns a literacy three-way of $-0.0350$ ($t = -2.97$) and a college three-way of $+0.00076$ ($t = +0.65$); literacy is not a proxy for education. Third, the within-stock IO contrast discriminates against any generic state-level moderator that hits all headquarters-state stocks alike. Fourth, the size-controlled specification discriminates against the most natural firm-level moderator. What remains is a class of ownership-interactive state-level confounds --- a state-level moderator that itself operates differently across IO terciles --- which the design does not separate; the canonical instance is state-resident public-pension-fund ownership, which is concentrated in roughly seven states with active state-pension systems and is collinear with the state fixed effect for the remaining forty-plus states. Section~\ref{sec:discussion} engages this limitation.

\paragraph{Correct headquarters-state assignment.} The static Compustat snapshot mis-assigns relocator firm-months for the 14 percent of observations whose static state differs from the firm's then-current 10-K headquarters state. Headline patterns are estimated on the stable-HQ subsample, for which the static and point-in-time merges coincide and the local-investor unit of analysis is well-defined. The PIT literacy reassignment on the full panel is the strictest available correction; the resulting magnitude collapse (Section~\ref{sec:robustness}) is the binding population-scope limitation of the result.

\paragraph{Cluster-robust inference with 51 state clusters.} The natural clustering dimension is the state, and with 51 clusters asymptotic cluster-robust standard errors over-reject. The binding inference procedure is the one-way state-level wild-cluster bootstrap of \citet{cameronGelbachMiller2008} with Webb six-point weights \citep{webb2014} and the null imposed. The two-way state-by-month CGM statistic is reported as a diagnostic in the Internet Appendix. For the cross-state literacy gradient, the precision-weighted-slope $t$-statistic is interpreted alongside a state-permutation $p$-value with $B = 9{,}999$.

\subsection{Causal language}

Under the conditional-unconfoundedness assumption, $\hat\gamma$ in Equation~\eqref{eq:size_ctrl} recovers the average within-state-month differential sensitivity of returns to $\mathrm{mom}\times\mathrm{IV}$ as a function of state-aggregate literacy, partialing out state and month means and size moderation of the same anomaly. It is the conditional association the design supports; it is not a clean average treatment effect of literacy.

I attempted an instrumental-variables strategy using state K-12 personal-finance graduation requirements as $Z$ for state literacy, building on the Urban-Schmeiser 2015/2020 database. The instrument fails its first-stage relevance test for documented reasons (mandate adoption was a state-level response to low adult literacy, producing a wrong-signed cross-state correlation between $Z$ and the moderator). A cohort-matched implementation using NFCS micro-data with state, wave, and age-band fixed effects produces a first-stage $F = 0.05$. The 2SLS estimates are directionally consistent with OLS but with standard errors inflated by a factor of 6--7. The OLS estimate from Equation~\eqref{eq:size_ctrl} is the preferred inferential statement throughout the paper. Section~\ref{sec:robustness} documents the IV failure in detail.
